\begin{tabularx}{\textwidth}[t]{|X|X|}
  \hline 
  \textbf{Notions} & \textbf{Commentaires} \\
  \hline 
  Jeux d'accessibilité à deux joueurs sur un graphe. Stratégie. Stratégie gagnante. Position gagnante. Détermination des positions gagnantes par le calcul des attracteurs. Construction de stratégies gagnantes.
   & 
   On considère des jeux à deux joueurs $\left(J_1\right.$ et $\left.J_2\right)$ modélisés par des graphes bipartis (l'ensemble des états contrôlés par $J_1$ et l'ensemble des états contrôlés par $J_2$ ). Il y a trois types d'états finals : les états gagnants pour $J_1$, les états gagnants pour $J_2$ et les états de match nul.
On ne considère que les stratégies sans mémoire.
  \\
  \hline 
  Notion d'heuristique. Algorithme minmax avec une heuristique.
   & 
   L'élagage alpha-beta n'est pas au programme.
  \\
  \hline 
  \multicolumn{2}{|X|}{\textbf{Mise en œuvre}}
  \\
  \hline 
  \multicolumn{2}{|p{0.95\textwidth}|}{
  La connaissance dans le détail des algorithmes de cette section n'est pas un attendu du programme. Les étudiants acquièrent une familiarité avec les idées sous-jacentes qu'ils peuvent réinvestir dans des situations où les modélisations et les recommandations d'implémentation sont guidées, notamment dans leurs aspects arborescents.
  } 
  \\
  \hline
\end{tabularx}
